\section{Introdução}
\label{sec:introducao}

O desenvolvimento de veículos aéreos não tripulados (VANTs) envolve a integração de diversos subsistemas responsáveis pelo funcionamento adequado da aeronave, incluindo estruturas mecânicas, sistemas de propulsão, controle de voo e sistemas eletrônicos embarcados. Nesse contexto, um setor de Sistemas Embarcados desempenha um papel fundamental na integração e no gerenciamento dos componentes eletrônicos responsáveis na operação de um projeto de um drone autônomo. \\
Os sistemas embarcados são responsáveis pela interface entre sensores, atuadores e unidades de processamento, garantindo que informações críticas de navegação e controle sejam adquiridas, processadas e transmitidas de forma confiável durante todas as fases de operação do veículo. Além disso, este setor também atua diretamente no planejamento da arquitetura elétrica do sistema, na distribuição de energia para os diferentes módulos eletrônicos e na integração entre hardware e software utilizados no drone. No contexto da Delta V, o setor de Sistemas Embarcados tem como objetivo assegurar que todos os componentes eletrônicos do Logan, nosso principal drone, operem de forma integrada, estável e segura, contribuindo para a confiabilidade do sistema e para o sucesso das missões propostas na competição EletroQuad SAE BRASIL – AXIA Energia 2026.\\ 
Este relatório apresenta as atividades desenvolvidas pelo setor ao longo do ciclo de desenvolvimento da aeronave para a competição de 2026, abordando a arquitetura do sistema embarcado, os critérios utilizados na seleção dos componentes eletrônicos, a disposição dos módulos na aeronave e os testes realizados para validação do sistema. Além disso, busca registrar as informações e ações do setor, assegurando que o conhecimento adquirido não se perca com o tempo, transformando-o em uma referência sólida e uma base de estudos para futuros membros da equipe. Essa documentação não apenas preserva a memória técnica do setor, mas também otimiza projetos futuros, permitindo que as próximas gerações avancem a partir dos aprendizados e soluções já desenvolvidas. \\\

\subsection{Escopo do setor}
\subsubsection{Objetivos do setor}

\begin{itemize}
    \item Desenvolver e documentar a arquitetura elétrica e de integração do sistema embarcado;
    \item Garantir comunicação confiável entre os módulos e distribuição estável de energia;
    \item Otimizar a disposição física dos componentes para reduzir ruídos e interferências;
    \item Implementar soldagem e montagem seguras, assegurando a integridade do sistema.
\end{itemize}

\subsubsection{Responsabilidades do setor}

\begin{itemize}
    \item Esquematização completa do sistema elétrico do drone, definindo conexões e interligações entre controladores, módulos de alimentação e atuadores;
    \item Modelagem em ferramentas CAD para prototipagem, simulação e elaboração do desenho de circuitos, garantindo precisão no desenvolvimento;
    \item Soldagem e montagem de cabos, fios e componentes, utilizando técnicas apropriadas à estação de retrabalho e ferro de solda, assegurando durabilidade das conexões;
    \item Definição estratégica da disposição dos componentes para otimizar espaço interno e compatibilidade eletromagnética (EMC);
    \item Análise de logs resultantes de voos bem ou mal sucedidos, atuando como uma caixa-preta com o fim de retornar para os demais setores, informações como autonomia do drone, potência dos motores e esforço de controle, dados úteis para a programação do VANT, por exemplo;
    \item Execução de testes, simulações e validação da integração entre subsistemas elétricos e de controle, garantindo confiabilidade no funcionamento geral do veículo.
\end{itemize}